% =====================================================================================================
% Section 5: Colour registers.  Body only; to be \input by the main file.
% Requires: amsmath, amssymb, amsthm, booktabs; environments theorem, lemma, proposition, corollary,
% definition, remark (one shared counter), proof; macros \F, \C, \N, \R; \cite{oai};
% external labels thm:grouped, lem:batch, lem:certificate, prop:interface.
% =====================================================================================================
\section{Colour registers}\label{sec:colour}

\newcommand{\colBot}{\mathrm{Bot}}
\newcommand{\colTop}{\mathrm{Top}}
\newcommand{\colXb}{\mathrm{Xb}}
\newcommand{\colYt}{\mathrm{Yt}}
\newcommand{\colDD}{\mathbb{D}}
\newcommand{\colK}{\widetilde K}

In the network of \cite[\S2]{oai}, every invocation carries one side wire $A_{ST}$ for each ordered pair
of neighbouring triples. That is $d=\binom{h-3}{3}+3(h-3)$ auxiliary wires per bank element, and these
wires make up almost all of the role count $W$. In this section we replace most of them by
\emph{registers}. Each bank element owns a few auxiliary wires. The registers of many source elements
meet the registers of many target elements at a common binary label, where a single gate can apply an
arbitrary transfer matrix between them. A twelve-step schedule realizes the required transfer for
arbitrary auxiliary contents (Lemma~\ref{col:lem:scalar}), and the stage-2 invocation is its
inverse-transpose (Lemma~\ref{col:lem:reverse}). For the stage sizes $(h_1,h_2,h_3)=(30,20,41)$ the
network uses about $57$, $40$ and $73$ auxiliary wires per bank element, instead of $3006$, $731$ and
$8550$. It satisfies the hypotheses of the grouped tensor-power theorem with exponent
$z=1-3.21\cdot10^{-6}$ (Theorem~\ref{col:thm:main}).

We use the notation of \cite[\S2]{oai} throughout and cite its results without re-proving them. The
construction is in Sections~\ref{col:sec:data}--\ref{col:sec:network}: factor data and centres, meeting
labels, register designs, the invocations, labels and edges, and the network. The designs follow in
Section~\ref{col:sec:designs}, and the exponent and the main theorem in
Sections~\ref{col:sec:exponent}--\ref{col:sec:main}. Each statement below is proved here, cited, or
explicitly marked as verified by exhaustive computation. The whole argument of this section is part of
the formal verification described in Section~\ref{sec:formal}.

%-----------------------------------------------------------------------------------------------------
\subsection{Factor data and centres}\label{col:sec:data}

Fix integers $h_1,h_2,h_3\ge 9$, indexed by the stage order. For $j=1,2,3$ let $D_j=\F^{h_j}$ with its
ordinary dot product, and let
\[
  \colDD=D_1\otimes D_2\otimes D_3=\F^{m},\qquad m=h_1h_2h_3,
\]
with the tensor dot product, which is the ordinary dot product of $\F^m$ \cite[\S2.3]{oai}. For a
binary vector $z$ we have $z\cdot z\equiv\mathrm{wt}(z)\pmod 2$, so $z$ has norm one if and only if it
has odd weight. A subspace is \emph{nondegenerate} if the dot product restricted to it has zero
radical. An \emph{orthonormal basis} (ONB) is a basis $z_1,\dots,z_r$ with $z_i\cdot z_k=\delta_{ik}$.
Orthogonal complements are taken in the ambient space indicated ($D_j$, a tensor factor, or $\colDD$).
For nested subspaces $U\subseteq V$, the \emph{residual} is $V\cap U^\perp$ \cite[Lemma~2.2]{oai}. We
use two standard facts about symmetric bilinear forms over $\F$.
\begin{itemize}
\item[(F1)] If $U\subseteq V$ and $U$ is nondegenerate, then $V=U\perp(V\cap U^\perp)$ (an orthogonal direct
  sum), so $\dim(V\cap U^\perp)=\dim V-\dim U$. The residual $V\cap U^\perp$ is nondegenerate if and only if
  $V$ is. If $V=U\perp W$ is any orthogonal decomposition with $U$ nondegenerate, then $V\cap U^\perp=W$.
  In a nondegenerate ambient space of dimension $n$, $\dim U^\perp=n-\dim U$.
\item[(F2)] The Gram matrix of a tensor product of bases is the Kronecker product of the Gram
  matrices. Hence tensor products of nondegenerate subspaces are nondegenerate, and tensor products
  of ONBs are ONBs. An orthogonal sum is nondegenerate if and only if its summands are, and the union
  of ONBs of the summands is an ONB of the sum.
\end{itemize}

\begin{definition}[factor data]\label{col:def:factor}
Fix $j$ and write $h=h_j$, $D=D_j$. Let $\mathcal T=\mathcal T_j$ be the set of $3$-subsets of
$[h]=\{1,\dots,h\}$, $v=v_j=\binom h3$, and $t_T\in D$ the indicator of $T\in\mathcal T$, so
$t_T\cdot t_T=1$. As in \cite[\S2.1]{oai}, $S,T\in\mathcal T$ are \emph{neighbours} if $|S\cap T|$ is
even, equivalently $t_S\cdot t_T=0$; neighbours are distinct. An \emph{ordered orthogonal pair} is an
ordered pair $(T,S)$ of neighbours, so $|T\cap S|\in\{0,2\}$. Let $\Omega=\Omega_j$ be the set of
ordered orthogonal pairs, $\Omega^0$ the subset with $T\cap S=\emptyset$, and $\Omega^2$ the subset with
$|T\cap S|=2$. Then $\Omega=\Omega^0\sqcup\Omega^2$, with
$|\Omega^0|=v\binom{h-3}{3}$ and $|\Omega^2|=3v(h-3)$.
\end{definition}

The banks are those of \cite[\S2.1]{oai}: wires $X_d,Y_d$ for
$d\in\mathcal T_1\times\mathcal T_2\times\mathcal T_3$, with $N_0=v_1v_2v_3$, $u_d=t_{d_1}\otimes
t_{d_2}\otimes t_{d_3}$ and $U_d=\langle u_d\rangle$. At stage $j$, an \emph{invocation} is indexed by
$o\in\mathcal O_j:=\prod_{l\neq j}\mathcal T_l$, so there are $N_0/v_j$ invocations. For
$T\in\mathcal T_j$ let $(o,T)$ be the key with $j$-th entry $T$ and other entries $o$. The invocation
$o$ has \emph{source bank} $x_T:=X_{(o,T)}$ and \emph{target bank} $y_S:=Y_{(o,S)}$
($T,S\in\mathcal T_j$). Its stage spaces are those of \cite[\S2.3]{oai}:
\[
  E=\textstyle\bigotimes_{l<j}D_l,\quad P=\langle p\rangle,\ p=\bigotimes_{l<j}t_{o_l},\quad B=P^\perp\cap E,
  \quad F=\bigotimes_{l>j}D_l,\quad Q=\langle q\rangle,\ q=\bigotimes_{l>j}t_{o_l}.
\]
An empty tensor product is $\F$. Thus $p\cdot p=q\cdot q=1$ and $E=B\perp P$, with $B=0$ at stage~1
and $Q=F=\F$ at stage~3. Put $\alpha_j=\dim E=\prod_{l<j}h_l$ and $\beta_j=\dim F=\prod_{l>j}h_l$. Then
$m=\alpha_jh_j\beta_j$ and $\dim B=\alpha_j-1$.

\begin{definition}[centres]\label{col:def:centres}
With $h$, $\mathcal T$ as above, let $G\colon\C^{\mathcal T}\to\C^{h}$ and $R\colon\C^{h}\to\C^{\mathcal T}$ be
\[
  (Gx)_i=\tfrac12\Big(\sum_{T\ni i}x_T-\tfrac13\sum_{T\in\mathcal T}x_T\Big),\qquad
  (Rc)_S=\sum_{i\in S}c_i ,
\]
and for $S,T\in\mathcal T$ put $K_{ST}:=\tfrac12\big(1-|S\cap T|\big)$. Thus $K_{ST}=K_{TS}$; it equals
$\tfrac12$ on disjoint pairs and $-\tfrac12$ on pairs meeting in two points.
\end{definition}

\begin{lemma}[$h$ centres]\label{col:lem:centres}
For all $S,T\in\mathcal T$,
\[
  (RG)_{ST}=\tfrac12\big(|S\cap T|-1\big)=-K_{ST}.
\]
Consequently $RG+K\circ 1_\Omega=I$, where $(K\circ1_\Omega)_{ST}=K_{ST}$ if $(T,S)\in\Omega$ and $0$ otherwise.
\end{lemma}

\begin{proof}
Since $|S|=3$,
$(RG)_{ST}=\sum_{i\in S}G_{iT}=\tfrac12\sum_{i\in S}\big([i\in T]-\tfrac13\big)=\tfrac12(|S\cap T|-1)$.
Consider the entry $(S,T)$ of $RG+K\circ1_\Omega$. If $S=T$, then $|S\cap T|=3$, so $(RG)_{TT}=1$, and
$(T,T)\notin\Omega$. If $|S\cap T|=1$, then $(RG)_{ST}=0$ and $(T,S)\notin\Omega$. If $|S\cap T|\in\{0,2\}$,
then $(T,S)\in\Omega$ and the entry is $(RG)_{ST}+K_{ST}=0$.
\end{proof}

\begin{remark}\label{col:rem:centres}
OpenAI's network uses $h+1$ centres $c_1,\dots,c_h,c_*$ and obtains the same matrix
$(RG)_{ST}=(|S\cap T|-1)/2$ \cite[\S2.1 and proof of Lemma~2.1]{oai}. Lemma~\ref{col:lem:centres} shows
that $c_*$ is unnecessary: each invocation needs only $h_j$ centres. OpenAI's side-wire map is
$(JA)_S=-\sum_{T}\frac{|S\cap T|-1}{2}A_{ST}=\sum_{T}K_{ST}A_{ST}$, the sum over the neighbours $T$
of $S$.
\end{remark}

%-----------------------------------------------------------------------------------------------------
\subsection{Meeting labels}\label{col:sec:labels}

Fix a stage $j$ and an invocation $o$, and write $D=D_j$, $h=h_j$, $\alpha=\alpha_j$, $\beta=\beta_j$.

\begin{definition}[meeting labels]\label{col:def:lambda}
For a subspace $L\subseteq D$ put
\[
  \Lambda_L:=\big((B\otimes D)\perp(P\otimes L)\big)\otimes Q\ \subseteq\ \colDD .
\]
The special labels are
\[
  \colBot:=\Lambda_0=B\otimes D\otimes Q,\qquad \colXb(t):=\Lambda_{\langle t\rangle},\qquad
  \colYt(t):=\Lambda_{t^\perp},\qquad \colTop:=\Lambda_D=E\otimes D\otimes Q,
\]
for $t=t_T$, with $t^\perp$ taken in $D$.
\end{definition}

These are labels of \cite[(2.7)]{oai}, with $D$ replaced by $D_j$. $\colBot$ is the row-1 and row-4
label. $\colXb(t)$ is the row-2 label. $\colYt(t)$ is the row-5 label and the outgoing label of the
target bank. $\colTop$ is the row-3, row-6 and row-7 label and the outgoing label of the source bank.
The row-0 label $B\otimes\langle t\rangle\otimes Q$ of \cite{oai} is replaced by $\colBot$.

\begin{lemma}[meeting labels]\label{col:lem:lambda}
Let $L\subseteq L'\subseteq D$ be subspaces.
\begin{enumerate}
\item[(i)] $\Lambda_L=\colBot\perp(P\otimes L\otimes Q)$, and $z\mapsto p\otimes z\otimes q$ is an
  isometric linear embedding of $D$ onto $P\otimes D\otimes Q$. Hence $\Lambda_L$ is nondegenerate if and
  only if $L$ is, and $\Lambda_L\subseteq\Lambda_{L'}$. In particular $\colBot$, $\colXb(t)$, $\colYt(t)$
  and $\colTop$ are nondegenerate.
\item[(ii)] If $L$ and $L'$ are nondegenerate, the residual of $\Lambda_L\subseteq\Lambda_{L'}$ is
  $P\otimes(L'\cap L^\perp)\otimes Q$, of dimension $\dim L'-\dim L$.
\item[(iii)] If $z_1,\dots,z_r$ is an ONB of $N\subseteq D$, then $p\otimes z_1\otimes q,\dots,p\otimes z_r\otimes q$
  is an ONB of $P\otimes N\otimes Q$.
\item[(iv)] \cite[end of \S2.3]{oai} A nondegenerate subspace $N\subseteq\F^n$ has an ONB if and only if
  $N=0$ or $N$ contains a vector of odd weight.
\end{enumerate}
\end{lemma}

\begin{proof}
(i) $B\perp P$ gives $B\otimes D\perp P\otimes L$. Moreover
$(p\otimes z\otimes q)\cdot(p\otimes z'\otimes q)=(p\cdot p)(z\cdot z')(q\cdot q)=z\cdot z'$, and the map is
injective because $p,q\ne0$. So $P\otimes L\otimes Q$ is isometric to $L$. The space $B$ is
nondegenerate by (F1), being the residual of the nondegenerate line $P$ in $E$. Hence $\colBot$ is
nondegenerate by (F2), and $\Lambda_L$ is nondegenerate if and only if $L$ is, again by (F2). The
subspaces $0$, $\langle t\rangle$, $t^\perp$ and $D$ are nondegenerate, $t^\perp$ by (F1). Monotonicity is
clear.
(ii) By (F1), $L'=L\perp N$ with $N=L'\cap L^\perp$. Applying the isometry of (i),
$\Lambda_{L'}=\Lambda_L\perp(P\otimes N\otimes Q)$, and the last part of (F1) identifies the residual. Its
dimension is $\dim N=\dim L'-\dim L$. (iii) follows from the isometry in (i). (iv) is the splitting
statement at the end of \cite[\S2.3]{oai}, with ``norm one'' read as ``odd weight''. The converse is
trivial, since ONB vectors have norm one.
\end{proof}

%-----------------------------------------------------------------------------------------------------
\subsection{Register designs}\label{col:sec:designs-def}

\begin{definition}[register design]\label{col:def:design}
Let $h\ge 9$, $D=\F^h$, and let $\mathcal T$, $\Omega$ be as in Definition~\ref{col:def:factor}. A
\emph{register design} $\mathcal R$ for $D$ consists of finite sets $\mathcal X$ (X-registers),
$\mathcal Y$ (Y-registers), $\mathrm{Tr}$ (transfers) and $\Sigma$ (side pairs) with the following
data.
\begin{itemize}
\item Each $\xi\in\mathcal X$ has an \emph{owner} $T_\xi\in\mathcal T$ and a \emph{chain}
  $L_{\xi,1}\subsetneq\dots\subsetneq L_{\xi,k_\xi}$ of subspaces of $D$, with $k_\xi\ge1$.
\item Each $\eta\in\mathcal Y$ has an owner $S_\eta\in\mathcal T$ and a chain
  $L_{\eta,1}\subsetneq\dots\subsetneq L_{\eta,k_\eta}$, with $k_\eta\ge1$.
\item Each $\tau\in\mathrm{Tr}$ has an X-register $\xi_\tau$, a Y-register $\eta_\tau$ and a label
  $L_\tau\subseteq D$. Its \emph{owner pair} is $(T_{\xi_\tau},S_{\eta_\tau})$.
\item Each $\sigma\in\Sigma$ is an ordered orthogonal pair $(T_\sigma,S_\sigma)\in\Omega$.
\end{itemize}
The \emph{meeting labels} are the subspaces that occur in chains. For $T,S\in\mathcal T$ the
\emph{cover count} is
\[
  n(T,S)=\#\{\tau:(T_{\xi_\tau},S_{\eta_\tau})=(T,S)\}+\#\{\sigma:(T_\sigma,S_\sigma)=(T,S)\}.
\]
The design must satisfy:
\begin{enumerate}
\item[(D1)] every meeting label is nondegenerate;
\item[(D2)] for every X-register with owner $T$ and chain $L_1\subsetneq\dots\subsetneq L_k$:
  $t_T\in L_1$, and each of the spaces $L_1\cap t_T^\perp$, $L_{i+1}\cap L_i^\perp$ $(1\le i<k)$ and
  $L_k^\perp$ is $0$ or contains a vector of odd weight;
\item[(D3)] for every Y-register with owner $S$ and chain $L_1\subsetneq\dots\subsetneq L_k$:
  $L_k\subseteq t_S^\perp$, and each of the spaces $L_1$, $L_{i+1}\cap L_i^\perp$ $(1\le i<k)$ and
  $t_S^\perp\cap L_k^\perp$ is $0$ or contains a vector of odd weight;
\item[(D4)] for every transfer $\tau$, the label $L_\tau$ belongs to the chain of $\xi_\tau$ and to the
  chain of $\eta_\tau$;
\item[(D5)] $n(T,S)=1$ for $(T,S)\in\Omega$, and $n(T,S)=0$ for $(T,S)\notin\Omega$.
\end{enumerate}
For $\Omega'\subseteq\Omega$, a \emph{partial register design for $\Omega'$} satisfies (D1)--(D4) and
(D5$_{\Omega'}$): $n(T,S)=1$ for $(T,S)\in\Omega'$ and $n(T,S)=0$ otherwise. We write
$\mathrm{regs}(\mathcal R)=|\mathcal X|+|\mathcal Y|$ and $\mathrm{side}(\mathcal R)=|\Sigma|$.
\end{definition}

If $\mathcal X=\mathcal Y=\mathrm{Tr}=\emptyset$ and $\Sigma=\Omega$, (D1)--(D5) hold trivially. This
design gives OpenAI's side wires, one for each ordered orthogonal pair. The point of registers is that
one meeting label can serve many pairs: if $w$ X-registers and $r$ Y-registers share a meeting label,
transfers at that label can cover up to $wr$ ordered pairs with $w+r$ wires.

\begin{lemma}[union]\label{col:lem:union}
If $\mathcal R'$ and $\mathcal R''$ are partial register designs for disjoint sets $\Omega',\Omega''$, then
their disjoint union (of registers, transfers and side pairs) is a partial register design for
$\Omega'\cup\Omega''$. In particular, a partial design for $\Omega^0$ together with one for $\Omega^2$ is a
register design.
\end{lemma}

\begin{proof}
Conditions (D1)--(D4) concern individual registers and transfers, and the cover counts add.
\end{proof}

\begin{lemma}[first-cover designs]\label{col:lem:firstcover}
Let $\Omega'\subseteq\Omega$, and let $\mathcal X$, $\mathcal Y$ be finite sets of registers (owners and
chains) satisfying (D1)--(D3). Let $\mathcal C$ be a finite, totally ordered set of \emph{candidates}
$\gamma=(L_\gamma,\xi_\gamma,\eta_\gamma)$, where $L_\gamma$ lies in the chains of both $\xi_\gamma$ and
$\eta_\gamma$ and the owner pair $(T_{\xi_\gamma},S_{\eta_\gamma})$ lies in $\Omega'$. Call a pair
\emph{covered} if it is the owner pair of some candidate. Define:
\begin{itemize}
\item $\mathrm{Tr}$: for each covered pair, the first candidate with that owner pair;
\item $\Sigma$: the pairs of $\Omega'$ that are not covered.
\end{itemize}
Then $(\mathcal X,\mathcal Y,\mathrm{Tr},\Sigma)$ is a partial register design for $\Omega'$, and
$|\Sigma|=|\Omega'|-\#\{\text{covered pairs}\}$.
\end{lemma}

\begin{proof}
(D1)--(D3) are assumed, and (D4) holds because every candidate's label lies in both chains. A pair of
$\Omega'$ is covered by exactly one transfer if it is covered, and by exactly one side pair otherwise.
A pair outside $\Omega'$ is the owner pair of no candidate and is not in $\Sigma$.
\end{proof}

\begin{remark}[first-cover semantics]\label{col:rem:firstcover}
In Lemma~\ref{col:lem:firstcover}, a candidate that is not selected still has both registers passing
through its label. In step~5 of the schedule below the two registers then meet at that label with
coefficient $0$, which is harmless. A design can therefore be specified by its registers and an ordered
candidate list. The transfer of a pair is assigned to its first covering candidate, and the side pairs
are \emph{defined} as the uncovered pairs. With these semantics, the exact cover (D5) holds by
construction. The only global quantity that must be computed is the number $|\Sigma|$ of uncovered
pairs, which enters the role count.
\end{remark}

%-----------------------------------------------------------------------------------------------------
\subsection{The invocations}\label{col:sec:invocations}

Fix a stage $j$, an invocation $o$, and data $(\mathcal X,\mathcal Y,\mathrm{Tr},\Sigma)$ as in
Definition~\ref{col:def:design} for $D=D_j$. Conditions (D1)--(D5) are imposed only where stated. The
wires of the invocation are:
\begin{itemize}
\item the source and target bank entries $x_T$, $y_S$ $(T,S\in\mathcal T)$;
\item centres $c_1,\dots,c_h$;
\item one side wire $A_\sigma$ for each $\sigma\in\Sigma$;
\item one X-register wire $a_\xi$ for each $\xi\in\mathcal X$ and one Y-register wire $b_\eta$ for each
  $\eta\in\mathcal Y$.
\end{itemize}
All but the bank entries are private to the invocation. Besides $G$ and $R$ we use the linear maps
\begin{gather*}
  (\Pi x)_\xi=x_{T_\xi},\qquad (Vx)_\sigma=x_{T_\sigma},\qquad
  (\Xi b)_S=\sum_{\eta:\,S_\eta=S}b_\eta,\qquad (JA)_S=\sum_{\sigma:\,S_\sigma=S}K_{ST_\sigma}A_\sigma,\\
  (\colK_La)_\eta=\sum_{\tau:\,\eta_\tau=\eta,\ L_\tau=L}K_{S_\eta T_{\xi_\tau}}\,a_{\xi_\tau},\qquad
  \colK=\sum_{L}\colK_L .
\end{gather*}
So $\Pi$ copies $x_T$ into the X-registers of $T$, $\Xi$ sums the Y-registers of $S$ into $y_S$, and
$\colK$ is the \emph{transfer matrix}: the entry $K_{ST}$ goes from the X-register to the Y-register of
each transfer with owner pair $(T,S)$. When $\Sigma=\Omega$, the map $J$ is OpenAI's
(Remark~\ref{col:rem:centres}).

\begin{definition}[forward invocation]\label{col:def:forward}
The \emph{forward invocation} is the chronological sequence of gates in
Table~\ref{col:tab:schedule}. A gate is a linear map $I+Z$ of the vector of wire values. $Z$ is the
displayed increment, and it is supported on the listed wires, which the gate \emph{touches}. The gate's
label is the listed subspace. In step~5 there is one gate for every meeting label $L$, with increment
$\colK_L$. It touches all registers whose chain contains $L$, and the gates are applied in order of
nondecreasing $\dim L$, with ties in any fixed order. Under (D4), $\colK_L$ acts only on registers whose
chains contain $L$, so its support lies in the touched wires.
\end{definition}

\begin{table}[ht]\centering\small
\begin{tabular}{@{}r l p{6.9cm} l@{}}
\toprule
step & update & gates and touched wires & label\\
\midrule
1 & $y\leftarrow y-Rc$ & one gate: all $y_S$ and all $c_i$ & $\colBot$\\
2 & $y\leftarrow y-\Xi b-JA$ & one gate per $S$: $y_S$; $b_\eta$ with $S_\eta=S$; $A_\sigma$ with $S_\sigma=S$ & $\colBot$\\
3 & $b\leftarrow b-\colK a$ & one gate: all registers & $\colBot$\\
4 & $a\leftarrow a+\Pi x$, $A\leftarrow A+Vx$ & one gate per $T$: $x_T$; $a_\xi$ with $T_\xi=T$; $A_\sigma$ with $T_\sigma=T$ & $\colXb(t_T)$\\
5 & $b\leftarrow b+\colK_L a$ & one gate per meeting label $L$, by nondecreasing $\dim L$: all registers whose chain contains $L$ & $\Lambda_L$\\
6 & $c\leftarrow c+Gx$ & one gate: all $x_T$ and all $c_i$ & $\colTop$\\
7 & $y\leftarrow y+Rc$ & one gate: all $y_S$ and all $c_i$ & $\colBot$\\
8 & $y\leftarrow y+\Xi b+JA$ & one gate per $S$, touching the same wires as in step 2 & $\colYt(t_S)$\\
9 & $c\leftarrow c-Gx$ & as in step 6 & $\colTop$\\
10 & $b\leftarrow b-\colK a$ & as in step 3 & $\colTop$\\
11 & $a\leftarrow a-\Pi x$, $A\leftarrow A-Vx$ & one gate per $T$, touching the same wires as in step 4 & $\colTop$\\
12 & $b\leftarrow b+\colK a$ & as in step 3 & $\colTop$\\
\bottomrule
\end{tabular}
\caption{The forward invocation (stages 1 and 3). All labels include the factor $Q$
(Definition~\ref{col:def:lambda}).}\label{col:tab:schedule}
\end{table}

Gates of one step touch pairwise disjoint sets of wires, except in step~5. There, two gates at distinct
meeting labels of equal dimension touch disjoint sets of registers, because a strictly increasing chain
contains at most one subspace of each dimension. Every increment $Z$ maps wires of one kind into wires
of a different kind, as in
$x\to c$, $c\to y$, $\{b,A\}\to y$, $x\to\{a,A\}$ and $a\to b$. Hence $Z^2=0$ for every gate.

\begin{remark}\label{col:rem:oairows}
If $\mathcal X=\mathcal Y=\emptyset$, steps 3, 5, 10 and 12 are empty. The remaining steps are OpenAI's
rows \cite[(2.3)]{oai} in the order $1,0,2,3,4,5,6,7$, with $h$ centres and the row-0 label $\colBot$.
Rows~0 and~1 commute, since both only subtract from $y$, using wires that neither modifies.
\end{remark}

\begin{lemma}[scalar identity]\label{col:lem:scalar}
For arbitrary initial values of all wires, the forward invocation maps
\[
  (x,y,c,A,a,b)\longmapsto(x,\;y+Mx,\;c,A,a,b),\qquad M=RG+JV+\Xi\colK\Pi ,
\]
where $M_{ST}=(RG)_{ST}+K_{ST}\,n(T,S)$. If (D5) holds, then $M=I$. That is, the forward invocation
maps $(x,y)\mapsto(x,y+x)$ and restores every auxiliary wire.
\end{lemma}

\begin{proof}
No gate writes $x$. Let $(x,y_0,c_0,A_0,a_0,b_0)$ be the initial values. In step~5 every gate writes
only $b$ and reads only $a$, and $a$ is not written in step~5. So step~5 as a whole adds
$\sum_L\colK_La=\colK a$ to $b$, whatever the order of its gates. Steps with several gates act on
disjoint wires, so each step performs the displayed update. Table~\ref{col:tab:trace} records the
wires written in each step and their values afterwards. Every entry follows from the previous ones and
the update of the step.

\begin{table}[ht]\centering\small
\begin{tabular}{@{}r l l@{}}
\toprule
after step & wires written & value\\
\midrule
1 & $y$ & $y_0-Rc_0$\\
2 & $y$ & $y_0-Rc_0-\Xi b_0-JA_0$\\
3 & $b$ & $b_0-\colK a_0$\\
4 & $a$, $A$ & $a_0+\Pi x$,\quad $A_0+Vx$\\
5 & $b$ & $b_0-\colK a_0+\colK(a_0+\Pi x)=b_0+\colK\Pi x$\\
6 & $c$ & $c_0+Gx$\\
7 & $y$ & $y_0-Rc_0-\Xi b_0-JA_0+R(c_0+Gx)=y_0+RGx-\Xi b_0-JA_0$\\
8 & $y$ & $y_0+RGx-\Xi b_0-JA_0+\Xi(b_0+\colK\Pi x)+J(A_0+Vx)=y_0+Mx$\\
9 & $c$ & $c_0$\\
10 & $b$ & $b_0+\colK\Pi x-\colK(a_0+\Pi x)=b_0-\colK a_0$\\
11 & $a$, $A$ & $a_0$,\quad $A_0$\\
12 & $b$ & $b_0-\colK a_0+\colK a_0=b_0$\\
\bottomrule
\end{tabular}
\caption{Values in the proof of Lemma~\ref{col:lem:scalar}.}\label{col:tab:trace}
\end{table}

Now compute $M$. By definition,
$(JVx)_S=\sum_{\sigma:\,S_\sigma=S}K_{ST_\sigma}x_{T_\sigma}$, so
$(JV)_{ST}=K_{ST}\cdot\#\{\sigma:(T_\sigma,S_\sigma)=(T,S)\}$. Also
$(\Xi\colK\Pi x)_S=\sum_{\eta:\,S_\eta=S}\sum_{\tau:\,\eta_\tau=\eta}K_{ST_{\xi_\tau}}x_{T_{\xi_\tau}}$, so
$(\Xi\colK\Pi)_{ST}=K_{ST}\cdot\#\{\tau:(T_{\xi_\tau},S_{\eta_\tau})=(T,S)\}$. Hence
$M_{ST}=(RG)_{ST}+K_{ST}\,n(T,S)$. Under (D5), $n(T,S)=[(T,S)\in\Omega]$, and
Lemma~\ref{col:lem:centres} gives $M=I$.
\end{proof}

\begin{lemma}[reverse invocation]\label{col:lem:reverse}
Define the \emph{reverse invocation} by replacing each gate $I+Z$ of the forward invocation by
$I-Z^{\mathsf T}$, keeping its touched wires, its label and its place in the chronological order. Assume
(D5). For arbitrary initial values of all wires the reverse invocation maps
\[
  (x,y,c,A,a,b)\longmapsto(x-y,\;y,\;c,A,a,b).
\]
Every wire has the same sequence of touches and labels in both invocations, so the two invocations have
the same edges and residuals.
\end{lemma}

\begin{proof}
Let $N_1,\dots,N_r$ be the gates of the forward invocation in chronological order. By
Lemma~\ref{col:lem:scalar} and (D5), $N:=N_r\cdots N_1=I+\mathcal E$, where $\mathcal E$ has entry $1$ in
position $(y_T,x_T)$ for each $T$ and $0$ elsewhere. Since $Z^2=0$ for every gate,
$(I+Z)^{-1}=I-Z$, so the reverse gates are $N_i^{-\mathsf T}$. The reverse invocation is therefore
\[
  N_r^{-\mathsf T}\cdots N_1^{-\mathsf T}=\big(N_1^{-1}\cdots N_r^{-1}\big)^{\mathsf T}=(N^{-1})^{\mathsf T}
  =(I-\mathcal E)^{\mathsf T}=I-\mathcal E^{\mathsf T},
\]
using $\mathcal E^2=0$. The map $I-\mathcal E^{\mathsf T}$ sends $x_T\mapsto x_T-y_T$ and fixes every other wire.
The support of $Z^{\mathsf T}$ lies in the same set of wires as the support of $Z$, so the touched wires
and labels are unchanged by construction.
\end{proof}

\begin{remark}\label{col:rem:reverse}
Explicitly, the reverse invocation reads, for example: step~1 $c\leftarrow c+R^{\mathsf T}y$;
step~2 $b_\eta\leftarrow b_\eta+y_{S_\eta}$ and $A_\sigma\leftarrow A_\sigma+K_{S_\sigma T_\sigma}y_{S_\sigma}$;
step~4 $x_T\leftarrow x_T-\sum_{T_\xi=T}a_\xi-\sum_{T_\sigma=T}A_\sigma$; step~5
$a\leftarrow a-\colK_L^{\mathsf T}b$; step~6 $x\leftarrow x-G^{\mathsf T}c$. Data flow from $Y$ to $X$,
but every X-register still belongs to the X bank and every Y-register to the Y bank, so the physical
touch sequence is that of Table~\ref{col:tab:schedule}. This replaces OpenAI's reversal of the rows
\cite[\S2.3, after (2.7)]{oai} and makes a separate reverse schedule unnecessary.
\end{remark}

%-----------------------------------------------------------------------------------------------------
\subsection{Labels and edges}\label{col:sec:edges}

At stage $j$ the \emph{incoming labels} are $E\otimes\langle t_T\rangle\otimes Q$ on $x_T$ and
$B\otimes\langle t_S\rangle\otimes Q$ on $y_S$. The \emph{outgoing labels} are $\colTop$ on $x_T$ and
$\colYt(t_S)$ on $y_S$. Each auxiliary wire has source label $0$ and sink label $\colDD$. These are the
stage-boundary labels of \cite[\S2.3]{oai}, with $D$ replaced by $D_j$. As shown there, the outgoing
labels of stage $j$ are the incoming labels of stage $j+1$; the incoming labels at stage 1 are $U_d$
(on $X_d$) and $0$ (on $Y_d$); and the outgoing labels at stage 3 are $\colDD$ and $U_d^\perp$.

\begin{lemma}[labels and edges]\label{col:lem:edges}
Assume (D1)--(D4) for the design of factor $j$. In every invocation at stage $j$, forward or reverse:
\begin{enumerate}
\item[(a)] every gate label is nondegenerate;
\item[(b)] the label sequence of every wire is as in Table~\ref{col:tab:paths}. In particular each
  register is touched at $\colBot$, then (X-registers only) at $\colXb(t_T)$, then at its chain labels in
  increasing order, then (Y-registers only) at $\colYt(t_S)$, and finally at $\colTop$;
\item[(c)] consecutive labels along every wire are nested. The nonzero residuals and their dimensions are
  those of Table~\ref{col:tab:residuals}. Every residual is nondegenerate and has an ONB;
\item[(d)] every edge is increasing, except the centre edges from step~6 to step~7, which decrease by
  $P\otimes D\otimes Q$, of dimension $h_j$.
\end{enumerate}
\end{lemma}

\begin{table}[ht]\centering\small
\begin{tabular}{@{}l l@{}}
\toprule
wire & label sequence (step numbers in brackets)\\
\midrule
$x_T$ & in $\to\colXb(t_T)\,[4]\to\colTop\,[6,9,11]=$ out\\
$y_S$ & in $\to\colBot\,[1,2,7]\to\colYt(t_S)\,[8]=$ out\\
$c_i$ & $0\to\colBot\,[1]\to\colTop\,[6]\to\colBot\,[7]\to\colTop\,[9]\to\colDD$\\
$A_\sigma$ & $0\to\colBot\,[2]\to\colXb(t_{T_\sigma})\,[4]\to\colYt(t_{S_\sigma})\,[8]\to\colTop\,[11]\to\colDD$\\
$a_\xi$ & $0\to\colBot\,[3]\to\colXb(t_{T_\xi})\,[4]\to\Lambda_{L_{\xi,1}}\to\dots\to\Lambda_{L_{\xi,k}}\,[5]\to\colTop\,[10,11,12]\to\colDD$\\
$b_\eta$ & $0\to\colBot\,[2,3]\to\Lambda_{L_{\eta,1}}\to\dots\to\Lambda_{L_{\eta,k}}\,[5]\to\colYt(t_{S_\eta})\,[8]\to\colTop\,[10,12]\to\colDD$\\
\bottomrule
\end{tabular}
\caption{Label sequences in one invocation. ``In'' and ``out'' are the stage-boundary
labels.}\label{col:tab:paths}
\end{table}

\begin{table}[ht]\centering\small
\begin{tabular}{@{}l l l l@{}}
\toprule
wire & edge & residual & dimension\\
\midrule
$x_T$, $t=t_T$ & in $\to4$ & $B\otimes t^\perp\otimes Q$ & $(\alpha-1)(h-1)$\\
 & $4\to6$ & $P\otimes t^\perp\otimes Q$ & $h-1$\\
$y_S$, $t=t_S$ & in $\to1$ & $B\otimes t^\perp\otimes Q$ & $(\alpha-1)(h-1)$\\
 & $7\to8$ & $P\otimes t^\perp\otimes Q$ & $h-1$\\
$c_i$ & source $\to1$ & $B\otimes D\otimes Q$ & $(\alpha-1)h$\\
 & $1\to6$;\ $6\to7$ (decrease);\ $7\to9$ & $P\otimes D\otimes Q$ & $h$ each\\
 & $9\to$ sink & $E\otimes D\otimes Q^\perp$ & $\alpha h(\beta-1)$\\
$A_\sigma$, $(T,S)=(T_\sigma,S_\sigma)$ & source $\to2$ & $B\otimes D\otimes Q$ & $(\alpha-1)h$\\
 & $2\to4$ & $P\otimes\langle t_T\rangle\otimes Q$ & $1$\\
 & $4\to8$ & $P\otimes\langle t_T,t_S\rangle^\perp\otimes Q$ & $h-2$\\
 & $8\to11$ & $P\otimes\langle t_S\rangle\otimes Q$ & $1$\\
 & $11\to$ sink & $E\otimes D\otimes Q^\perp$ & $\alpha h(\beta-1)$\\
$a_\xi$, owner $T$ & source $\to3$ & $B\otimes D\otimes Q$ & $(\alpha-1)h$\\
\quad chain $L_1\subsetneq\dots\subsetneq L_k$ & $3\to4$ & $P\otimes\langle t_T\rangle\otimes Q$ & $1$\\
 & $4\to L_1$ & $P\otimes(L_1\cap t_T^\perp)\otimes Q$ & $\dim L_1-1$\\
 & $L_i\to L_{i+1}$ & $P\otimes(L_{i+1}\cap L_i^\perp)\otimes Q$ & $\dim L_{i+1}-\dim L_i$\\
 & $L_k\to10$ & $P\otimes L_k^\perp\otimes Q$ & $h-\dim L_k$\\
 & $12\to$ sink & $E\otimes D\otimes Q^\perp$ & $\alpha h(\beta-1)$\\
$b_\eta$, owner $S$ & source $\to2$ & $B\otimes D\otimes Q$ & $(\alpha-1)h$\\
\quad chain $L_1\subsetneq\dots\subsetneq L_k$ & $3\to L_1$ & $P\otimes L_1\otimes Q$ & $\dim L_1$\\
 & $L_i\to L_{i+1}$ & $P\otimes(L_{i+1}\cap L_i^\perp)\otimes Q$ & $\dim L_{i+1}-\dim L_i$\\
 & $L_k\to8$ & $P\otimes(t_S^\perp\cap L_k^\perp)\otimes Q$ & $h-1-\dim L_k$\\
 & $8\to10$ & $P\otimes\langle t_S\rangle\otimes Q$ & $1$\\
 & $12\to$ sink & $E\otimes D\otimes Q^\perp$ & $\alpha h(\beta-1)$\\
\bottomrule
\end{tabular}
\caption{Residuals of one invocation at stage $j$ ($h=h_j$, $\alpha=\alpha_j$, $\beta=\beta_j$).
``$L_i$'' denotes the step-5 touch at $\Lambda_{L_i}$, $Q^\perp$ is taken in $F$, and $\perp$ inside the
middle factor is taken in $D$. All edges not listed have zero residual.}\label{col:tab:residuals}
\end{table}

\begin{proof}
(a) The gate labels are $\colBot$, $\colTop$, $\colXb(t)$, $\colYt(t)$ and the $\Lambda_L$ for meeting
labels $L$. They are nondegenerate by Lemma~\ref{col:lem:lambda}(i) and (D1).

(b) Read off the touched wires of each step in Table~\ref{col:tab:schedule}. In step~5 the gates are
applied in order of nondecreasing dimension, and the dimensions along a chain strictly increase. So each
register meets its chain labels in chain order, and by (D4) each transfer lies in a gate touching both
of its registers. By Lemma~\ref{col:lem:reverse} the reverse invocation has the same sequences.

(c) We use the orthogonal decompositions $E=B\perp P$, $D=\langle t\rangle\perp t^\perp$ and
$F=Q\perp Q^\perp$, together with (F1) and Lemma~\ref{col:lem:lambda}(ii).
\begin{itemize}
\item \emph{Banks.} $E\otimes\langle t\rangle\otimes Q=(B\otimes\langle t\rangle\perp P\otimes\langle t\rangle)\otimes Q$ and
  $\colXb(t)=(B\otimes\langle t\rangle\perp B\otimes t^\perp\perp P\otimes\langle t\rangle)\otimes Q$. Hence the
  first is contained in the second, with residual $B\otimes t^\perp\otimes Q$. In the same way
  $B\otimes\langle t\rangle\otimes Q\subseteq\colBot$, with the same residual. The edges
  $\colXb(t)\subseteq\colTop$ and $\colBot\subseteq\colYt(t)$ are the cases $\langle t\rangle\subseteq D$
  and $0\subseteq t^\perp$ of Lemma~\ref{col:lem:lambda}(ii).
\item \emph{Centres.} $\colTop=\colBot\perp P\otimes D\otimes Q$ by Lemma~\ref{col:lem:lambda}(i), and
  $\colDD=E\otimes D\otimes F=\colTop\perp E\otimes D\otimes Q^\perp$.
\item \emph{Side wires.} $(T,S)\in\Omega$, so $\langle t_T\rangle\subseteq t_S^\perp$, and
  Lemma~\ref{col:lem:lambda}(ii) applies to $0\subseteq\langle t_T\rangle\subseteq t_S^\perp\subseteq D$.
  The residuals are $\langle t_T\rangle$, then $t_S^\perp\cap t_T^\perp=\langle t_T,t_S\rangle^\perp$, then
  $\langle t_S\rangle$.
\item \emph{Registers.} By (D2), $\langle t_T\rangle\subseteq L_1\subsetneq\dots\subsetneq L_k\subseteq D$. By
  (D3), $0\subseteq L_1\subsetneq\dots\subsetneq L_k\subseteq t_S^\perp\subseteq D$. All of these subspaces are
  nondegenerate by (D1) and (F1). Lemma~\ref{col:lem:lambda}(ii) gives the listed residuals and
  dimensions, using $\dim t_S^\perp=h-1$.
\end{itemize}
Every residual is nondegenerate by (F1). It remains to find ONBs. By (F2) and
Lemma~\ref{col:lem:lambda}(iii), it suffices to find ONBs of the tensor factors.
\begin{itemize}
\item $D$ and $E$ have the coordinate ONBs, and $\langle t\rangle$ has the ONB $t$.
\item $t^\perp$ contains $e_i$ for $i\notin T$. Since $(T,S)\in\Omega$, $\langle t_T,t_S\rangle^\perp$
  contains $e_i$ for $i\notin T\cup S$; such $i$ exists because $|T\cup S|\le6<h$.
\item If $B\ne0$, it contains a coordinate unit $e_{i_1}\otimes\dots\otimes e_{i_{j-1}}$ of $E$ with
  $i_1\notin o_1$, which is orthogonal to $p$. Similarly $Q^\perp\ne0$ contains a coordinate unit of
  $F$ orthogonal to $q$ \cite[proof of Lemma~2.3]{oai}.
\item The register residuals $L_1\cap t_T^\perp$, $L_{i+1}\cap L_i^\perp$, $L_k^\perp$ (X) and $L_1$,
  $L_{i+1}\cap L_i^\perp$, $t_S^\perp\cap L_k^\perp$ (Y) are $0$ or contain a vector of odd weight, by
  (D2) and (D3).
\end{itemize}
In every case Lemma~\ref{col:lem:lambda}(iv) gives an ONB.

(d) Along each sequence of Table~\ref{col:tab:paths} the labels increase, except the centres'
$\colTop\to\colBot$ between steps~6 and~7.
\end{proof}

\begin{remark}\label{col:rem:edges}
For banks, centres and side wires, the labels are those of \cite[(2.7)]{oai}, with $D_j$ in place of $D$,
the row-0 label $\colBot$ and $h$ centres. If the design has no registers, Lemma~\ref{col:lem:edges}
reduces to \cite[Lemma~2.3]{oai} for this variant. Conditions (D1)--(D3) can equivalently be
expressed by an explicit ONB of $D_j$ adapted to each register's chain. An X-register satisfies (D1) and (D2) if and only if $D$ has an ONB that contains $t_T$ and in which every
$L_i$ is spanned by a subset of the basis containing $t_T$. One direction takes the union of ONBs of
the summands of $D=\langle t_T\rangle\perp(L_1\cap t_T^\perp)\perp(L_2\cap L_1^\perp)\perp\dots\perp L_k^\perp$,
which exist by (F1) and Lemma~\ref{col:lem:lambda}(iv). Conversely, every residual is then spanned by
basis vectors. The same holds for Y-registers, with $t_S$ in the basis but outside every $L_i$.
\end{remark}

%-----------------------------------------------------------------------------------------------------
\subsection{The network}\label{col:sec:network}

\begin{definition}[colour-register network]\label{col:def:network}
Let $h_1,h_2,h_3\ge9$ and let $\mathcal R_j$ be a register design for $D_j$ ($j=1,2,3$). The network
$\mathcal N=\mathcal N(\mathcal R_1,\mathcal R_2,\mathcal R_3)$ has the bank wires $X_d,Y_d$. For every
stage $j$ and every invocation $o\in\mathcal O_j$ it also has a private copy of the auxiliary wires of
Section~\ref{col:sec:invocations} for $\mathcal R_j$: $h_j$ centres, $\mathrm{side}(\mathcal R_j)$ side wires and
$\mathrm{regs}(\mathcal R_j)$ register wires. The stages are run in the order $1,2,3$. Stage $j$ consists of
the invocations $o\in\mathcal O_j$, with $x_T=X_{(o,T)}$ and $y_S=Y_{(o,S)}$. They are forward
invocations at stages 1 and 3 and reverse invocations at stage 2, and their labels are formed from the
stage spaces of $o$. The terminal labels are those of \cite[\S2.3]{oai}. Every source has label $0$,
except $X_d$, which has label $U_d$. Every sink has label $\colDD$, except $Y_d$, which has label
$U_d^\perp$.
\end{definition}

The invocations of one stage act on pairwise disjoint sets of wires. Indeed $X_d$ and $Y_d$ belong
only to the invocation $o=(d_l)_{l\ne j}$, and the auxiliary wires are private. So the order of the
invocations within a stage is immaterial.

\begin{proposition}[the network]\label{col:prop:network}
Let $\mathcal R_1,\mathcal R_2,\mathcal R_3$ be register designs and $\mathcal N=\mathcal N(\mathcal R_1,\mathcal R_2,\mathcal R_3)$.
\begin{enumerate}
\item[(a)] For arbitrary initial values on all wires, the three stages map
  $(X_d,Y_d)\mapsto(-Y_d,X_d)$ for every $d$ and restore every auxiliary wire.
\item[(b)] Every gate label is nondegenerate, all wires touched by a gate carry its label, and every edge
  has nested nondegenerate endpoint labels and a residual with an ONB. The only decreasing edges are the
  centre edges between steps~6 and~7.
\item[(c)] The number of wires and the saving $\Delta:=Wm-\sum_e\rho_e$ are
  \[
    W=2N_0+\sum_{j=1}^3\frac{N_0}{v_j}\big(h_j+\mathrm{side}(\mathcal R_j)+\mathrm{regs}(\mathcal R_j)\big),\qquad
    \Delta=2N_0-2\sum_{j=1}^3\frac{N_0}{v_j}h_j^2 ,
  \]
  where $\rho_e$ runs over the residual dimensions of all edges.
\item[(d)] The multiset of nonzero residual dimensions is the union over $j$ of $N_0/v_j$ copies of the
  stage-$j$ multiset of Table~\ref{col:tab:multiset}.
\end{enumerate}
Consequently $\mathcal N$ satisfies the network interface (\ref{prop:interface}), since the proof of
\cite[Proposition~2.4]{oai} applies verbatim. For every $f\ge1$, $\mathcal N$ applies
$C^{\otimes mf}$ to each of its $W$ arbitrary arrays, and every edge $e$ is one frame change of residual
dimension $\rho_e$ (Lemma~\ref{lem:batch}).
\end{proposition}

\begin{table}[ht]\centering\small
\begin{tabular}{@{}l l l@{}}
\toprule
class & wires per invocation & nonzero edge dimensions of each wire\\
\midrule
banks & $2v_j$ & $(\alpha-1)(h-1)$,\ $h-1$\\
centres & $h_j$ & $(\alpha-1)h$,\ $h$,\ $h$,\ $h$,\ $\alpha h(\beta-1)$\\
side wires & $\mathrm{side}(\mathcal R_j)$ & $(\alpha-1)h$,\ $1$,\ $h-2$,\ $1$,\ $\alpha h(\beta-1)$\\
X-registers & $|\mathcal X_j|$ & $(\alpha-1)h$,\ $1$,\ $d_1-1$,\ $d_2-d_1,\dots,d_k-d_{k-1}$,\ $h-d_k$,\ $\alpha h(\beta-1)$\\
Y-registers & $|\mathcal Y_j|$ & $(\alpha-1)h$,\ $d_1$,\ $d_2-d_1,\dots,d_k-d_{k-1}$,\ $h-1-d_k$,\ $1$,\ $\alpha h(\beta-1)$\\
\bottomrule
\end{tabular}
\caption{Edge multiset of one invocation at stage $j$ ($h=h_j$, $\alpha=\alpha_j$, $\beta=\beta_j$; zero
entries omitted). For a register, $d_1<\dots<d_k$ are the dimensions of its chain labels. The bank edges
are those entering stage $j$ and those inside it.}\label{col:tab:multiset}
\end{table}

\begin{proof}
(a) By Lemmas~\ref{col:lem:scalar} and~\ref{col:lem:reverse}, and because the invocations of a stage act
on disjoint wires, stage~1 maps $Y_d\mapsto Y_d+X_d$, stage~2 maps $X_d\mapsto X_d-Y_d$, and stage~3
maps $Y_d\mapsto Y_d+X_d$. Each stage restores every auxiliary wire. Hence
$(x,y)\mapsto(x,y+x)\mapsto(-y,y+x)\mapsto(-y,x)$, as in \cite[Lemma~2.1]{oai}.

(b) Within an invocation this is Lemma~\ref{col:lem:edges}. Every auxiliary wire lies in exactly one
invocation, where it runs from its source label $0$ to its sink label $\colDD$. A bank wire passes through
the three stages. Its source and sink labels are the incoming labels of stage~1 and the outgoing labels
of stage~3, and the outgoing labels of stage $j$ are the incoming labels of stage $j+1$
(Section~\ref{col:sec:edges}). So its edges are exactly the bank edges of Lemma~\ref{col:lem:edges} for
the three stages.

(c) The formula for $W$ counts the wires. For $\Delta$ we follow \cite[proof of Lemma~2.3]{oai}. Along
every wire, the signed changes of label dimension telescope to (sink dimension) $-$ (source dimension).
The source dimensions sum to $N_0$ and the sink dimensions to $Wm-N_0$, so the signed sum is
$Wm-2N_0$. By (b), the only decreases are one edge of dimension $h_j$ on each of the
$(N_0/v_j)h_j$ centres of stage $j$. Passing to absolute values adds twice these losses:
$\sum_e\rho_e=Wm-2N_0+2\sum_j(N_0/v_j)h_j^2$.

(d) Read off Table~\ref{col:tab:residuals}, using $\dim B=\alpha-1$, $\dim E=\alpha$ and
$\dim Q^\perp=\beta-1$. Each bank wire lies in exactly one invocation per stage.

For the consequence: the proof of \cite[Proposition~2.4]{oai} uses three things. The first is the
complete linear map \cite[(2.6)]{oai}, whose scalar part is the signed bank exchange of (a). The second
is the label properties of (b), through \cite[Lemma~2.2]{oai}. The third is the terminal labels, through
the odd weight $\mathrm{wt}(u_d)=27$. All three hold for $\mathcal N$.
\end{proof}

%-----------------------------------------------------------------------------------------------------
\subsection{Designs}\label{col:sec:designs}

We construct partial designs for $\Omega^2$ by a structured construction and for $\Omega^0$ from
searched data. Throughout, $h\ge9$ and $D=\F^h$.

\subsubsection*{Sunflower labels: pairs meeting in two points}

Fix an \emph{axis} $\{a,b\}\subseteq[h]$ and put $Q_{ab}=[h]\setminus\{a,b\}$, so $|Q_{ab}|=n:=h-2$. For
$x\in Q_{ab}$ let $g_x=e_a+e_b+e_x=t_{\{a,b,x\}}$, and for $P\subseteq Q_{ab}$ let
$M(P)=\mathrm{span}\{g_x:x\in P\}\subseteq D$.

\begin{lemma}[sunflower labels]\label{col:lem:sunlabels}
Let $P\subseteq Q_{ab}$.
\begin{enumerate}
\item[(a)] $g_x\cdot g_y=\delta_{xy}$. So $\{g_x\}_{x\in P}$ is an ONB of $M(P)$, $M(P)$ is nondegenerate
  of dimension $|P|$, and for $P\subseteq P'$ the residual of $M(P)\subseteq M(P')$ is $M(P'\setminus P)$.
\item[(b)] For $x\in P$: $t_{\{a,b,x\}}=g_x\in M(P)$ and $M(P)\cap g_x^\perp=M(P\setminus\{x\})$.
\item[(c)] If $P\ne Q_{ab}$, then $M(P)^\perp$ contains $e_y$ for every $y\in Q_{ab}\setminus P$.
\item[(d)] For $c\in Q_{ab}\setminus P$: $t_{\{a,b,c\}}\in M(P)^\perp$, and
  $t_{\{a,b,c\}}^\perp\cap M(P)^\perp=M(P\cup\{c\})^\perp$. This space contains a vector of odd weight if
  and only if $P\cup\{c\}\ne Q_{ab}$ or $h$ is even. If $P\cup\{c\}=Q_{ab}$ and $h$ is odd, it is nonzero and has
  no ONB.
\end{enumerate}
\end{lemma}

\begin{proof}
(a) $g_x\cdot g_y\equiv|\{a,b,x\}\cap\{a,b,y\}|\pmod 2$, which is $3$ or $2$. Orthonormal vectors are
independent, so they form an ONB of their span; the residual follows from (F1). (b) follows from (a).
(c) For $x\in P$ we have $e_y\cdot g_x=[y=x]=0$. (d) $g_c\perp M(P)$ by (a), and
$g_c^\perp\cap M(P)^\perp=(M(P)+\langle g_c\rangle)^\perp=M(P\cup\{c\})^\perp$. If $P\cup\{c\}\ne Q_{ab}$, use
(c). If $P\cup\{c\}=Q_{ab}$, then $M(Q_{ab})^\perp$ has dimension $h-(h-2)=2$ and contains the
independent vectors $w_a=e_a+1_{Q_{ab}}$ and $w_b=e_b+1_{Q_{ab}}$. Indeed, for $x\in Q_{ab}$ the only
nonzero contributions to $w_a\cdot g_x$ are $e_a\cdot e_a=1$ and $1_{Q_{ab}}\cdot e_x=1$, so
$w_a\cdot g_x=0$, and likewise for $w_b$. So
$M(Q_{ab})^\perp=\{0,w_a,w_b,e_a+e_b\}$, with weights $h-1$, $h-1$, $2$. It contains an odd vector if and
only if $h$ is even, and by Lemma~\ref{col:lem:lambda}(iv) it has no ONB for odd $h$.
\end{proof}

For every axis fix a cyclic order $x_0,\dots,x_{n-1}$ of $Q_{ab}$, with indices in $\mathbb Z/n$. For
$1\le\ell\le n$ let $U(i,\ell)=\{x_i,x_{i+1},\dots,x_{i+\ell-1}\}$, the arc of length $\ell$ starting at
$x_i$. Put
\[
  \ell_{\max}=n-1\ \text{ if $h$ is even},\qquad \ell_{\max}=n-2\ \text{ if $h$ is odd}.
\]

\begin{proposition}[cyclic sunflower design]\label{col:prop:sunflower}
The following data form a partial register design $\mathcal R^{\mathrm{sun}}_h$ for $\Omega^2$. For each axis
$A=\{a,b\}$ and $i,k\in\mathbb Z/n$:
\begin{itemize}
\item an X-register $\xi_{A,i}$ with owner $A\cup\{x_i\}$ and chain
  \[M(U(i,1))\subsetneq M(U(i,2))\subsetneq\dots\subsetneq M(U(i,\ell_{\max}));\]
\item a Y-register $\eta_{A,k}$ with owner $A\cup\{x_k\}$ and chain
  \[M(U(k-1,1))\subsetneq M(U(k-2,2))\subsetneq\dots\subsetneq M(U(k-\ell_{\max},\ell_{\max}));\]
\item for $i\ne k$ with $\ell:=(k-i)\bmod n\le\ell_{\max}$, a transfer
  $\tau_{A,i,k}=(M(U(i,\ell)),\xi_{A,i},\eta_{A,k})$;
\item if $h$ is odd, a side pair $\sigma_{A,k}=(A\cup\{x_{k+1}\},A\cup\{x_k\})$.
\end{itemize}
It has $3v$ X-registers, $3v$ Y-registers ($6$ per element), $3v\ell_{\max}$ transfers, and $3v$ side pairs
if $h$ is odd ($3$ per element), none if $h$ is even. Every chain has dimensions $1,2,\dots,\ell_{\max}$.
\end{proposition}

\begin{proof}
The arc $U(i,\ell)$ has $\ell$ elements, and $U(i,\ell+1)=U(i,\ell)\cup\{x_{i+\ell}\}$. The arc
$U(k-\ell,\ell)=\{x_{k-\ell},\dots,x_{k-1}\}$ does not contain $x_k$ for $\ell\le n-1$, and
$U(k-\ell-1,\ell+1)=U(k-\ell,\ell)\cup\{x_{k-\ell-1}\}$. Both chains are therefore strictly
increasing, with dimensions $1,\dots,\ell_{\max}$ by Lemma~\ref{col:lem:sunlabels}(a).

(D1) holds by Lemma~\ref{col:lem:sunlabels}(a).

(D2) Let $T=A\cup\{x_i\}$, so $t_T=g_{x_i}\in M(U(i,1))=\langle g_{x_i}\rangle$. The residuals are
$M(U(i,1))\cap t_T^\perp=0$ (part (b)); $M(\{x_{i+\ell}\})\ni g_{x_{i+\ell}}$, of weight $3$ (part (a));
and $M(U(i,\ell_{\max}))^\perp$, which contains some $e_y$ by part (c), since $\ell_{\max}<n$.

(D3) Let $S=A\cup\{x_k\}$, so $t_S=g_{x_k}$. Since $x_k\notin U(k-\ell,\ell)$, every chain label lies in
$t_S^\perp$ by part (a). The residuals $M(\{x_{k-1}\})$ and $M(\{x_{k-\ell-1}\})$ contain vectors of
weight $3$. By part (d) with $c=x_k$, the last residual is
$t_S^\perp\cap M(U(k-\ell_{\max},\ell_{\max}))^\perp=M(U(k-\ell_{\max},\ell_{\max}+1))^\perp$. The arc
$U(k-\ell_{\max},\ell_{\max}+1)$ equals $Q_{ab}$ exactly when $\ell_{\max}+1=n$, that is, when $h$ is even.
So part (d) gives an odd vector in both parities.

(D4) If $k\equiv i+\ell$, then $U(i,\ell)=U(k-\ell,\ell)$ is the $\ell$-th label of both chains.

(D5$_{\Omega^2}$) Let $(T,S)\in\Omega^2$. Then $A:=T\cap S$ is the unique axis contained in both, and
$T=A\cup\{x_i\}$, $S=A\cup\{x_k\}$ with $i\ne k$ uniquely determined. The owner pair of
$\tau_{A',i',k'}$ is $(A'\cup\{x_{i'}\},A'\cup\{x_{k'}\})$, and that of $\sigma_{A',k'}$ is
$(A'\cup\{x_{k'+1}\},A'\cup\{x_{k'}\})$. Both have intersection $A'$, so they equal $(T,S)$ only if
$A'=A$, and then only for $(i',k')=(i,k)$, respectively $k'=k$ and $i=k+1$. With
$\ell=(k-i)\bmod n\in\{1,\dots,n-1\}$ this gives $n(T,S)=[\ell\le\ell_{\max}]+[h\text{ odd},\ \ell=n-1]$.
For $h$ even, $\ell_{\max}=n-1$ and there are no side pairs. For $h$ odd, exactly one of $\ell\le n-2$
and $\ell=n-1$ holds. In both cases $n(T,S)=1$. Every owner pair meets in exactly two points, so pairs
outside $\Omega^2$ have $n=0$, and side pairs lie in $\Omega$.

Counts: there are $\binom h2$ axes with $n$ registers of each kind. Since $\binom h2(h-2)=3\binom h3$, there
are $3v$ registers of each kind. Each $i$ has $\ell_{\max}$ partners $k$, and for odd $h$ each $k$ gives one
side pair.
\end{proof}

\begin{remark}\label{col:rem:sunflower}
For odd $h$ the choice $\ell_{\max}=n-1$ would violate (D3). The last residual of the reader would be
$M(Q_{ab})^\perp$, which has no ONB by Lemma~\ref{col:lem:sunlabels}(d). The side pairs
$\sigma_{A,k}$ replace exactly those transfers.
\end{remark}

\subsubsection*{Coordinate labels: disjoint pairs}

For $A\subseteq[h]$ let $\F^A=\mathrm{span}\{e_i:i\in A\}\subseteq D$, and write $A^c=[h]\setminus A$.

\begin{lemma}[coordinate labels and chains]\label{col:lem:coord}
\begin{enumerate}
\item[(a)] $\F^A$ is nondegenerate with ONB $\{e_i\}_{i\in A}$, $(\F^A)^\perp=\F^{A^c}$, and for $A\subseteq A'$
  the residual of $\F^A\subseteq\F^{A'}$ is $\F^{A'\setminus A}$.
\item[(b)] An X-register with owner $T$ and chain $\F^{A_1}\subsetneq\dots\subsetneq\F^{A_k}$ satisfies (D1)
  and (D2) if $T\subseteq A_1$ and $|A_1|\ge4$.
\item[(c)] A Y-register with owner $S$ and chain $\F^{A_1}\subsetneq\dots\subsetneq\F^{A_k}$ satisfies (D1)
  and (D3) if $S\cap A_k=\emptyset$ and $|A_k|\le h-4$.
\item[(d)] Suppose $\F^{A}$ lies in the chain of an X-register as in (b), with owner $T$, and in the chain
  of a Y-register as in (c), with owner $S$. Then $T\cap S=\emptyset$; in particular $(T,S)\in\Omega^0$.
\end{enumerate}
\end{lemma}

\begin{proof}
(a) is clear.
(b) $t_T\in\F^{A_1}$ since $T\subseteq A_1$. The residual $\F^{A_1}\cap t_T^\perp$ contains $e_i$ for
$i\in A_1\setminus T$, which exists since $|A_1|>3$. The residual $\F^{A_{i+1}\setminus A_i}$ is
nonzero and contains a coordinate unit. Finally $(\F^{A_k})^\perp=\F^{A_k^c}$ is $0$ or contains a
coordinate unit.
(c) $\F^{A_k}\perp t_S$ since the supports are disjoint. The residuals $\F^{A_1}$ and
$\F^{A_{i+1}\setminus A_i}$ are $0$ or contain a coordinate unit. The last residual
$t_S^\perp\cap\F^{A_k^c}$ contains $e_x$ for $x\in A_k^c\setminus S$, which exists since $|A_k^c|\ge4>|S|$.
(d) Let $A_1$ be the first set of the X-chain and $A'$ the last set of the Y-chain. Then
$T\subseteq A_1\subseteq A\subseteq A'$ and $S\cap A'=\emptyset$, so $T\cap S=\emptyset$.
\end{proof}

The bounds are sharp. If $|A_1|=3$ then $A_1=T$, and $\F^T\cap t_T^\perp$ consists of the even-weight
vectors of $\F^T$, which has no ONB. The case $|A_k|=h-3$ is symmetric. A coordinate label $\F^A$ can
serve every writer $T\subseteq A$ and every reader $S\subseteq A^c$.

\begin{definition}[prefix-chain coordinate design]\label{col:def:prefix}
A \emph{prefix-chain coordinate design} for $D=\F^h$ is given by a list of groups $g=1,\dots,\Gamma$.
Group $g$ consists of the following data.
\begin{itemize}
\item A chain of sets $A_{g,1}\subsetneq\dots\subsetneq A_{g,k_g}\subseteq[h]$ with $4\le|A_{g,i}|\le h-4$ and
  $k_g\le9$.
\item \emph{Writer entries} $(T,g,i)$ with $T\subseteq A_{g,i}$.
\item \emph{Reader entries} $(S,g,I)$ with $\emptyset\ne I\subseteq\{1,\dots,k_g\}$ and
  $S\cap A_{g,\max I}=\emptyset$.
\end{itemize}
Its register design $\mathcal R^{\mathrm{coord}}$ has the following parts.
\begin{itemize}
\item An X-register for each writer entry, with owner $T$ and chain $(\F^{A_{g,i}})$, a single label.
\item A Y-register for each reader entry, with owner $S$ and chain $(\F^{A_{g,i}})_{i\in I}$ in increasing
  order.
\item Candidates $(\F^{A_{g,i}},\,(T,g,i),\,(S,g,I))$, one for each writer entry and each reader entry of
  the same group with $i\in I$. They are ordered lexicographically by $(g,i)$, with ties broken by a fixed
  order of the entries.
\item Transfers and side pairs given by first cover relative to $\Omega^0$
  (Lemma~\ref{col:lem:firstcover}).
\end{itemize}
In the designs below, $A_{g,i}=\pi_g(\{1,\dots,s_g+i-1\})$ for a permutation $\pi_g$ of $[h]$ and an offset
$s_g$. The sets are prefixes of a permutation, which explains the name. This structure is not used in
the proofs.
\end{definition}

\begin{lemma}\label{col:lem:prefix}
$\mathcal R^{\mathrm{coord}}$ is a partial register design for $\Omega^0$. Its X-registers have one label, its
Y-registers have at most $9$ labels, and
$\mathrm{side}(\mathcal R^{\mathrm{coord}})=|\Omega^0|-\#\{\text{covered pairs}\}$.
\end{lemma}

\begin{proof}
Lemma~\ref{col:lem:coord}(b) and (c) give (D1)--(D3), using $A_{g,i}\subseteq A_{g,\max I}$ for $i\in I$.
Lemma~\ref{col:lem:coord}(d) shows that every candidate's owner pair lies in $\Omega^0$. Now apply
Lemma~\ref{col:lem:firstcover}.
\end{proof}

\begin{proposition}[the three searched designs; computational verification]\label{col:prop:designs}
For $h\in\{30,20,41\}$ let $\mathcal R^{\mathrm{coord}}_h$ be the prefix-chain coordinate design given by the
searched data of Remark~\ref{col:rem:designs}. Let $\mathcal R_h=\mathcal R^{\mathrm{coord}}_h\sqcup\mathcal R^{\mathrm{sun}}_h$.
\begin{enumerate}
\item[(a)] The data satisfy the conditions of Definition~\ref{col:def:prefix}. The label sizes lie in
  $[11,19]$, $[6,14]$ and $[16,24]$ respectively, and the longest Y-chains have $9$, $8$ and $9$ labels.
\item[(b)] The numbers of writer entries, reader entries and uncovered disjoint pairs are as in
  Table~\ref{col:tab:designs}.
\item[(c)] $\mathcal R_h$ is a register design, satisfying (D1)--(D5), with the counts of
  Table~\ref{col:tab:designs}.
\end{enumerate}
Parts (a) and (b) are established by exhaustive computation, not by a proof by hand; see
Remark~\ref{col:rem:designs} for their formal verification. Part (c) follows
from (a), (b), Lemma~\ref{col:lem:prefix}, Proposition~\ref{col:prop:sunflower} and
Lemma~\ref{col:lem:union}.
\end{proposition}

\begin{table}[ht]\centering\small
\begin{tabular}{@{}l r r r@{}}
\toprule
 & $h=30$ (stage 1) & $h=20$ (stage 2) & $h=41$ (stage 3)\\
\midrule
$v=\binom h3$ & $4060$ & $1140$ & $10\,660$\\
invocations $N_0/v$ & $12\,152\,400$ & $43\,279\,600$ & $4\,628\,400$\\
ordered orthogonal pairs $|\Omega|$ & $12\,204\,360$ & $833\,340$ & $91\,143\,000$\\
\quad covered by coordinate transfers & $11\,860\,737$ & $772\,124$ & $89\,880\,593$\\
\quad covered by sunflower transfers $3v\ell_{\max}$ & $328\,860$ & $58\,140$ & $1\,183\,260$\\
\quad side pairs (coordinate + sunflower) & $14\,763+0$ & $3\,076+0$ & $47\,167+31\,980$\\
coordinate X-registers (writer entries) & $96\,881$ & $18\,090$ & $304\,414$\\
coordinate Y-registers (reader entries) & $94\,429$ & $17\,552$ & $330\,907$\\
sunflower registers $6v$ & $24\,360$ & $6\,840$ & $63\,960$\\
$\mathrm{regs}(\mathcal R_h)$ & $215\,670$ & $42\,482$ & $699\,281$\\
$\mathrm{side}(\mathcal R_h)$ & $14\,763$ & $3\,076$ & $79\,147$\\
auxiliary wires per invocation $h+\mathrm{side}+\mathrm{regs}$ & $230\,463$ & $45\,578$ & $778\,469$\\
per bank element $(\mathrm{side}+\mathrm{regs})/v$ & $56.8$ & $40.0$ & $73.0$\\
OpenAI's $d=\binom{h-3}3+3(h-3)$ & $3006$ & $731$ & $8550$\\
\bottomrule
\end{tabular}
\caption{The designs $\mathcal R_{30}$, $\mathcal R_{20}$, $\mathcal R_{41}$.}\label{col:tab:designs}
\end{table}

\begin{remark}[how Proposition~\ref{col:prop:designs} was verified]\label{col:rem:designs}
The designs were found by a heuristic search, which has no proof value. They were then checked
exhaustively, in exact integer arithmetic, by three independently written programs. Each program
checks the
conditions of (a) for every group and entry. It also checks that each of the $12\,204\,360$,
$833\,340$ and $91\,143\,000$ ordered orthogonal pairs is covered exactly once, by one transfer or by
one pair of the explicit side-pair list stored with the design. Equivalently, the stored list is the
set of uncovered pairs, and its length is the count in (b). The three programs agree on all counts in
Table~\ref{col:tab:designs}.

In the formal verification (Section~\ref{sec:formal}) the same group data define partial register
designs for $\Omega^0$, with first-cover semantics, and their validity is checked by the Lean kernel
itself. The count bounds are proved there: at most $304\,414$ X-registers, $330\,907$ Y-registers and
$47\,167$ side pairs at $h=41$, and likewise at $h=30,20$ with the numbers of
Table~\ref{col:tab:designs}. In the formal version each Y-register follows the maximal chain of its
group (every label of the group disjoint from its owner). Its edge multiset may therefore differ from
that of the searched data. It still satisfies the counts and the chain-length bounds used in
Proposition~\ref{col:prop:exponent}(d).
\end{remark}

%-----------------------------------------------------------------------------------------------------
\subsection{The exponent}\label{col:sec:exponent}

Let $\mathcal N=\mathcal N(\mathcal R_{30},\mathcal R_{20},\mathcal R_{41})$ for the stage order $(h_1,h_2,h_3)=(30,20,41)$. Then
$m=24\,600$, $N_0=49\,338\,744\,000$, and the stage parameters are as follows:
\begin{center}\small
\begin{tabular}{@{}c c c c c c@{}}
\toprule
stage $j$ & $h_j$ & $(\alpha_j,\beta_j)$ & $(\alpha_j-1)h_j$ & $\alpha_jh_j(\beta_j-1)$ & bank edges\\
\midrule
1 & 30 & $(1,820)$ & $0$ & $24\,570$ & $0$, $29$\\
2 & 20 & $(30,41)$ & $580$ & $24\,000$ & $551$, $19$\\
3 & 41 & $(600,1)$ & $24\,559$ & $0$ & $23\,960$, $40$\\
\bottomrule
\end{tabular}
\end{center}
Let $W^*=2^{43}$. The \emph{padded network} $\mathcal N^*$ adds $W^*-W$ padding roles. Each has the label
path $0\to\langle e\rangle\to\colDD$, for a coordinate vector $e$ of $\colDD$, is touched by no gate, and
has edges of dimension $1$ and $m-1$, with residual ONBs $\{e\}$ and the other coordinate vectors. Its
frame ratio is $\Phi_{\colDD}\Phi_0^{-1}=C^{\otimes m}$ \cite[\S2.2]{oai}. Let $\mathcal E$ be the multiset of
nonzero residual dimensions of all edges of $\mathcal N^*$, and
\[
  \Phi(z)=\sum_{\rho\in\mathcal E}\Big(\frac{\rho}{m}\Big)^{z}.
\]

\begin{proposition}[exponent certificate]\label{col:prop:exponent}
For the padded network $\mathcal N^*$:
\begin{enumerate}
\item[(a)] $W=8\,475\,019\,577\,600\le 2^{43}=8\,796\,093\,022\,208$, so there are $321\,073\,444\,608$
  padding roles. Moreover $\Delta=26\,618\,807\,200$, and $W^*m-\sum_{\rho\in\mathcal E}\rho=\Delta>0$.
\item[(b)] $1\le\rho\le m-1$ for every $\rho\in\mathcal E$, and $|\mathcal E|\ge2W^*$.
\item[(c)] $\Phi(1-\varepsilon_1)\le2^{43}$ for $\varepsilon_1=3.279039\cdot10^{-6}$ (exact edge multiset).
\item[(d)] $\Phi(1-\varepsilon_2)\le2^{43}$ for $\varepsilon_2=3.221860\cdot10^{-6}$, by an argument that uses
  only the counts of Table~\ref{col:tab:designs} and the chain lengths of Lemma~\ref{col:lem:prefix}.
\item[(e)] $\Phi(z)<2^{43}$ for $z=1-3.21\cdot10^{-6}$.
\end{enumerate}
\end{proposition}

\begin{proof}
(a) Substitute Table~\ref{col:tab:designs} into Proposition~\ref{col:prop:network}(c). Each padding role
contributes $1+(m-1)=m$ to $\sum\rho$, so $W^*m-\sum_{\rho\in\mathcal E}\rho=Wm-\sum_e\rho_e=\Delta$.

(b) By Table~\ref{col:tab:multiset} and the stage table above, the fixed sizes are at most $24\,570$, the
register middle edges at most $h_j-1$, and the padding edges $1$ and $m-1=24\,599$. Every role has at
least two nonzero edges. A bank wire has at least the stage-1 edge $29$ and the stage-2 edges $551$ and
$19$. A centre has three edges $h$. A side wire has two unit edges. An X-register has the unit edge
$\colBot\to\colXb$ and middle edges summing to $h-1$. A Y-register has the unit edge $\colYt\to\colTop$
and middle edges summing to $h-1$. A padding role has $1$ and $m-1$.

(c) Lemma~\ref{lem:certificate} bounds each term $(\rho/m)^{1-\varepsilon}$ by a rational number computed
from a rational upper bound for $\ln(m/\rho)$. The resulting bound for $\Phi(1-\varepsilon_1)$, evaluated in
exact rational arithmetic over the exact multiset $\mathcal E$ (38 distinct values of $\rho$, with the
chain dimensions of every register of $\mathcal R_h$), is at most $2^{43}$. This is a computation,
which re-runs the design checks of Remark~\ref{col:rem:designs} before evaluating the bound. In an
independent recomputation of the multiset, the root of $\Phi(p)=2^{43}$ was found to be $1-p=3.2791253\cdot10^{-6}\ge\varepsilon_1$.

(d) Put $\varphi(c)=(c/m)^z$ for $c\ge0$, with $0<z\le1$. Since $c\mapsto c^z$ is concave and increasing,
Jensen's inequality gives $\sum_{i=1}^{n}\varphi(c_i)\le n\,\varphi(S/n)$ whenever $c_i\ge0$ and $\sum c_i\le S$
(tangent bound). For integers $c\ge1$ also $\varphi(c)\le c\,\varphi(1)$ (crude bound). The \emph{middle
edges} of an X-register are its edges from the step-4 touch at $\colXb(t_T)$ to the step-10 touch at
$\colTop$. The middle edges of a Y-register are its edges from the step-3 touch at $\colBot$ to the
step-8 touch at $\colYt(t_S)$. A register with $k$ chain labels has $k+1$ middle edges, and by
Table~\ref{col:tab:residuals} their dimensions sum to $h-1$. Its other nonzero edges are
$(\alpha-1)h$, one unit edge and $\alpha h(\beta-1)$. By Lemma~\ref{col:lem:prefix}, a coordinate
X-register has $2$ middle edges and a coordinate Y-register at most $10$. Bounding
the coordinate registers by the tangent bound and the sunflower registers by the crude bound gives,
with $a=\alpha_j$, $b=\beta_j$ and $h=h_j$ at stage $j$,
\begin{align*}
  \Phi(z)\le{}&\sum_{j=1}^3\frac{N_0}{v_j}\Big[\,2v_j\,\varphi_{\mathrm{bk}}+h_j\,\varphi_{\mathrm{c}}
   +\mathrm{side}_j\,\varphi_{\mathrm{sd}}+X_j\,\varphi_{X}+Y_j\,\varphi_{Y}+6v_j\,\varphi_{\mathrm{sun}}\Big]\\
   &+(W^*-W)\big(\varphi(1)+\varphi(m-1)\big),\\
  \varphi_{\mathrm{bk}}={}&\varphi((a-1)(h-1))+\varphi(h-1),\qquad
  \varphi_{\mathrm{c}}=\varphi((a-1)h)+3\varphi(h)+\varphi(ah(b-1)),\\
  \varphi_{\mathrm{sd}}={}&\varphi((a-1)h)+2\varphi(1)+\varphi(h-2)+\varphi(ah(b-1)),\\
  \varphi_{X}={}&\varphi((a-1)h)+\varphi(1)+2\varphi(\tfrac{h-1}{2})+\varphi(ah(b-1)),\\
  \varphi_{Y}={}&\varphi((a-1)h)+\varphi(1)+10\,\varphi(\tfrac{h-1}{10})+\varphi(ah(b-1)),\\
  \varphi_{\mathrm{sun}}={}&\varphi((a-1)h)+h\,\varphi(1)+\varphi(ah(b-1)),
\end{align*}
where $\varphi(0)=0$, and $X_j$, $Y_j$, $\mathrm{side}_j$ are the coordinate X-registers, coordinate
Y-registers and side wires of Table~\ref{col:tab:designs}. Each $\varphi$-value at $z=1-\varepsilon_2$ is
bounded by Lemma~\ref{lem:certificate} with rational logarithm bounds, $22$ distinct arguments in all.
The resulting rational number is at most $2^{43}$ (exact arithmetic).

(e) Each term of $\Phi$ is strictly decreasing in $z$, since $0<\rho/m<1$. As
$1-3.21\cdot10^{-6}>1-\varepsilon_2$, part (e) follows from (d), and also from (c). The rational inequality
obtained directly from the bound of (d) at $z=1-3.21\cdot10^{-6}$ is part of the formal verification
(Section~\ref{sec:formal}).
\end{proof}

\begin{remark}\label{col:rem:exponent}
The argument of part (d) uses only the counts of Table~\ref{col:tab:designs} and the chain-length
bounds. It therefore applies to any designs whose per-factor counts and chain lengths are dominated by
these, and gives for them $W\le2^{43}$, $1\le\rho_e\le m-1$ for every edge, and
$\sum_e(\rho_e/m)^z<2^{43}$ at $z=1-3.21\cdot10^{-6}$.

For comparison, the per-direction rule of \cite[Theorem~2.6]{oai}, with the same packing into $2^{43}$
roles, gives only $1-\theta\approx\Delta/(W^*m\ln m)\approx1.2167\cdot10^{-8}$ for this network. Almost
all of $\sum\rho$ lies in edges with $\rho$ close to $m$, which the grouped recursion charges nearly
nothing.
\end{remark}

%-----------------------------------------------------------------------------------------------------
\subsection{The colour-register theorem}\label{col:sec:main}

\begin{theorem}[colour-register network]\label{col:thm:main}
Let $(h_1,h_2,h_3)=(30,20,41)$, let $\mathcal R_{30}$, $\mathcal R_{20}$, $\mathcal R_{41}$ be the register designs
of Proposition~\ref{col:prop:designs}, and let $\mathcal N^*$ be the network
$\mathcal N(\mathcal R_{30},\mathcal R_{20},\mathcal R_{41})$ padded to $W^*=2^{43}$ roles. Then $\mathcal N^*$ satisfies the
hypotheses of Theorem~\ref{thm:grouped} with $W^*=2^{43}$ and $m=24\,600$:
\begin{enumerate}
\item[(i)] $\mathcal N^*$ is a network on $W^*$ roles that satisfies the network interface
  (\ref{prop:interface}). For arbitrary contents of all roles and every $f\ge1$, it applies
  $C^{\otimes mf}$ to every role.
\item[(ii)] Every edge has nested nondegenerate endpoint labels, a residual with an ONB, and residual
  dimension $1\le\rho_e\le m-1$.
\item[(iii)] $\displaystyle\sum_e\Big(\frac{\rho_e}{m}\Big)^{z}<2^{43}$ for $z=1-3.21\cdot10^{-6}$.
\end{enumerate}
The network has $W=8\,475\,019\,577\,600$ non-padding roles and saving $\Delta=26\,618\,807\,200$.
Inequality (iii) also holds with $z=1-3.279039\cdot10^{-6}$, as $\le$. Consequently,
Theorem~\ref{thm:grouped} applies to $\mathcal N^*$ with exponent $z=1-3.21\cdot10^{-6}$.
\end{theorem}

\begin{proof}
By Proposition~\ref{col:prop:designs}, each $\mathcal R_h$ is a register design. Part (i) is
Proposition~\ref{col:prop:network} for $\mathcal N$. The padding roles are touched by no gate, and their frame
ratio is $C^{\otimes m}$. Part (ii) is Proposition~\ref{col:prop:network}(b) together with
Proposition~\ref{col:prop:exponent}(b) and the description of the padding roles. Part (iii) and the
remaining numbers are Proposition~\ref{col:prop:exponent}(a), (c) and (e).
\end{proof}

\begin{remark}[status of the proof]\label{col:rem:status}
The proof of Theorem~\ref{col:thm:main} has two kinds of input.
\begin{enumerate}
\item[(1)] \emph{Written proofs} in this section, for general parameters:
  Lemmas~\ref{col:lem:centres}, \ref{col:lem:lambda}, \ref{col:lem:union}, \ref{col:lem:firstcover},
  \ref{col:lem:scalar}, \ref{col:lem:reverse}, \ref{col:lem:edges}, \ref{col:lem:sunlabels},
  \ref{col:lem:coord}, \ref{col:lem:prefix} and Propositions~\ref{col:prop:network},
  \ref{col:prop:sunflower}. They rely on \cite[\S2]{oai} (Lemmas 2.1--2.3, Proposition~2.4 and the
  splitting statement at the end of \S2.3), on standard facts about symmetric bilinear forms over $\F$,
  and on Lemmas~\ref{lem:batch} and~\ref{lem:certificate}.
\item[(2)] \emph{Exhaustive computation} on the three searched coordinate designs
  (Proposition~\ref{col:prop:designs}(a),(b)), done by independent programs, and the exact rational
  evaluation of the certificate bounds (Proposition~\ref{col:prop:exponent}(c),(d)).
\end{enumerate}
In addition, Theorem~\ref{col:thm:main}, and with it Theorem~\ref{thm:main}, has been formally
verified (Section~\ref{sec:formal}). The formal proof covers the network for arbitrary register
designs, the three factor designs, whose data are checked by the Lean kernel, and the exponent
inequality in the form of Proposition~\ref{col:prop:exponent}(d) at $z=1-3.21\cdot10^{-6}$.
\end{remark}
